\documentclass[11pt, a4paper]{article}
\usepackage[utf8]{inputenc}
\usepackage[T1]{fontenc}
\usepackage[french]{babel}
\usepackage{geometry}
\usepackage{xcolor}
\usepackage{enumitem}
\usepackage{hyperref}

\geometry{hmargin=2.5cm, vmargin=2.5cm}

\definecolor{umonsRed}{RGB}{171,0,51}

\title{\color{umonsRed}\Huge \textbf{Guide de Soutenance \& Script Oral}\\ \large Classification de la Densité Mammaire BI-RADS (VinDr-Mammo)}
\author{\textbf{ALLA Abdelouahad} \\ Université de Mons (UMONS) -- CETIC}
\date{Juillet 2026}

\begin{document}

\maketitle
\tableofcontents
\newpage

% =============================================================================
\section{Slide 1 : Titre de la Présentation}
% =============================================================================
\begin{quote}
    \textit{« Bonjour Professeur, bonjour à tous. Je vais vous présenter le bilan et l'avancement de mon travail de mémoire sur la classification automatique de la densité mammaire selon la norme BI-RADS sur VinDr-Mammo, en utilisant des modèles hybrides de Deep Learning et une architecture d'Apprentissage Fédéré. »}
\end{quote}

% =============================================================================
\section{Slide 2 : Plan de la Présentation}
% =============================================================================
\begin{quote}
    \textit{« Mon exposé s'articule autour de 5 parties : le défi du déséquilibre de données, notre stratégie de Data Augmentation Ciblée, les résultats expérimentaux des modèles baselines et hybrides, la pipeline hiérarchique, et enfin nos travaux en cours et prochaines étapes. »}
\end{quote}

% =============================================================================
\section{Slide 3 : Le Défi Diagnostique de la Densité Mammaire}
% =============================================================================
\begin{quote}
    \textit{« En mammographie, les seins denses masquent les lésions tumorales car le tissu fibroglandulaire et les tumeurs apparaissent tous deux en blanc. Sur le dataset VinDr-Mammo (20 000 images), nous avons un déséquilibre massif de 1 pour 19.5 entre la classe majoritaire C (76 \%) et la classe rare A (<1 \%). Sans correction, un réseau prédit toujours la classe C et échoue totalement sur la classe A avec 0 \% de rappel. »}
\end{quote}

% =============================================================================
\section{Slide 4 : Division des Données (Patient-Split)}
% =============================================================================
\begin{quote}
    \textit{« Pour garantir une rigueur scientifique absolue, nous appliquons un découpage strict par patient (80 \% Train/Val soit 16 000 images, et 20 \% Test soit 4 000 images). Aucune image d'une même patiente ne se retrouve à la fois dans le train et dans le test, éliminant tout risque de 'data leakage'. »}
\end{quote}

% =============================================================================
\section{Slide 5 : Pourquoi les modèles échouaient sur la Classe A ?}
% =============================================================================
\begin{quote}
    \textit{« Au départ, le modèle ReXNet150 affichait 86.50 \% d'accuracy globale, mais avec 0 \% de rappel sur la classe A. Pour résoudre ce problème médicalement inacceptable sans déformer le jeu de validation, nous avons introduit une Data Augmentation Ciblée uniquement sur le train set, ramenant le ratio de 1:19.5 à 1:2.7. »}
\end{quote}

% =============================================================================
\section{Slide 6 : Mécanisme de Data Augmentation}
% =============================================================================
\begin{quote}
    \textit{« Concrètement, le réseau suréchantillonne les classes minoritaires A, B et D, et applique à la volée des transformations géométriques et photométriques (rotations, flips, variations de contraste). Résultat : le rappel sur la classe A est débloqué, passant de 0 \% à 87.5 \% de sensibilité ! »}
\end{quote}

% =============================================================================
\section{Slide 7 : Modèle ReXNet150 CC avec Data Augmentation (Figure 16)}
% =============================================================================
\begin{quote}
    \textit{« Sur la vue CC avec augmentation, ReXNet150 atteint 81.56 \% en validation et 80.00 \% en test final sur 2 000 images. Le rappel sur la classe critique D (très dense) monte à 87.0 \% en validation et 73.7 \% en test, démontrant un rééquilibrage réussi entre les 4 classes. »}
\end{quote}

% =============================================================================
\section{Slide 8 : Modèles Baselines Hybrides (ResNet50 \& ReXNet150 MLO)}
% =============================================================================
\begin{quote}
    \textit{« Sur les vues MLO, notre baseline ReXNet150 + Histogramme atteint notre meilleur score global d'accuracy à 86.50 \%. ResNet50 atteint 82.81 \%. Ces modèles profitent de la richesse de l'histogramme pour capturer la distribution globale des intensités. »}
\end{quote}

% =============================================================================
\section{Slide 9 : Ablation Study - Branches Image Seules (Section 6.2.4.3)}
% =============================================================================
\begin{quote}
    \textit{« Pour évaluer le rôle propre du réseau visuel sans l'histogramme, nous avons réalisé une étude d'ablation en entraînant les extracteurs visuels seuls sur 2 000 images de test. ResNet50 MLO atteint 83.00 \% d'accuracy de test et ReXNet150 CC atteint 82.90 \%. Cela confirme la solidité de nos extracteurs d'images. »}
\end{quote}

% =============================================================================
\section{Slide 10 : Tableau Récapitulatif Global des Expérimentations}
% =============================================================================
\begin{quote}
    \textit{« Ce tableau maître résume l'ensemble de nos 7 configurations. ReXNet150 MLO domine en accuracy globale (86.50 \%), tandis que ReXNet150 CC avec augmentation ciblée domine en équilibre des classes minoritaires B et D (80.5 \% et 73.7 \% de rappel). »}
\end{quote}

% =============================================================================
\section{Slide 11 : Approche Hiérarchique (Pipeline 3 Étapes)}
% =============================================================================
\begin{quote}
    \textit{« Nous avons aussi développé un pipeline hiérarchique à 3 étapes simulant la décision médicale : l'étape 1 sépare le non-dense (AB) du dense (CD). Pour les cas où l'étape 1 est confiante (>0.9), la précision de routage atteint 99.49 \% ! »}
\end{quote}

% =============================================================================
\section{Slide 12 : Architecture Hybride avec Texture GLCM (Section 6.2.5)}
% =============================================================================
\begin{quote}
    \textit{« Nous testons l'intégration de descripteurs de texture Haralick (GLCM : entropie, contraste, homogénéité). Les premiers entraînements atteignent 79.59 \% sur ReXNet150 CC et 75.50 \% sur ResNet50 MLO. »}
\end{quote}

% =============================================================================
\section{Slide 13 : État d'Avancement \& Travaux en Cours}
% =============================================================================
\begin{quote}
    \textit{« Actuellement, les modèles GLCM et ReXNet150 MLO finalisent leurs calculs sur le cluster GPU. La prochaine étape immédiate est le lancement de la Double Branche Image Multi-Vues (CC + MLO fusionnés en 1024D). Enfin, nous déploierons le modèle retenu dans la plateforme d'apprentissage fédéré Flower. »}
\end{quote}

% =============================================================================
\section{Slide 14 : Conclusion \& Points Forts du Bilan}
% =============================================================================
\begin{quote}
    \textit{« En conclusion, nous avons résolu le déséquilibre de classe sans fuite de données, validé la complémentarité des branches visuelles et textuelles, et obtenu des modèles atteignant jusqu'à 86.50 \% d'accuracy. Je vous remercie pour votre attention et je suis à votre disposition pour vos questions. »}
\end{quote}

\end{document}
